\documentclass[11pt,a4paper]{article}
\usepackage[margin=1in]{geometry}
\usepackage{amsmath,amssymb,amsthm}
\usepackage{algorithm}
\usepackage{algpseudocode}
\usepackage{booktabs}
\usepackage{hyperref}

\theoremstyle{definition}
\newtheorem{definition}{\textbf{Definition}}[section]
\newtheorem{theorem}{\textbf{Theorem}}[section]
\newtheorem{lemma}[theorem]{\textbf{Lemma}}
\newtheorem{remark}{\textbf{Remark}}[section]

\title{\textbf{Sequence-Level and On-Policy Distillation for Generative Language Models}}
\author{\textbf{Team Scaibu} \\ \small Email: \texttt{jaydeep.9664920749@gmail.com} \\ \small \textbf{Architecture and Core Engine Team}}
\date{\textbf{October 2026}}

\begin{document}
\maketitle

\section{Definitions and Cognitive Mental Models}

\subsection{Formal Mathematical Definition}
\textbf{Definition (Sequence and On-Policy Autoregressive Distillation):}
When distilling autoregressive language models, traditional word-level cross-entropy distillation on fixed static teacher transcripts suffers from \textbf{exposure bias}: during deployment, the student generates its own tokens step-by-step and inevitably visits off-distribution sequences where it was never trained to recover.

To resolve this, modern generative distillation encompasses three complementary paradigms:
\begin{enumerate}
    \item \textbf{Sequence-Level Distillation (Kim \& Rush, 2016)}:
    The teacher autoregressively generates complete candidate sequences $y^* \sim P_{\text{teacher}}(y \mid x)$ (e.g. via beam search or greedy decoding). The student is trained via standard Supervised Fine-Tuning (SFT) on the teacher's self-generated sequence tokens:
    {\boldmath
    \begin{equation}
    \mathcal{L}_{\text{seq-sft}} = - \sum_{t=1}^{|y^*|} \log P_{\text{student}}(y_t^* \mid x, y_{<t}^*)
    \end{equation}
    }
    \item \textbf{Token-Level Forward KL Distillation}:
    Matching next-token distributions across all vocabulary tokens $V$ using forward KL divergence:
    {\boldmath
    \begin{equation}
    \mathcal{L}_{\text{forward-KL}} = \sum_{t=1}^{|y|} D_{\text{KL}}\left( P_{\text{teacher}}(\cdot \mid x, y_{<t}) \,\|\, P_{\text{student}}(\cdot \mid x, y_{<t}) \right)
    \end{equation}
    }
    Forward KL mode is \emph{zero-avoiding} (mode-covering), forcing the student to place probability mass everywhere the teacher has mass, which can cause hallucination and blurriness.
    \item \textbf{On-Policy Reverse KL Distillation (MiniLLM, 2023)}:
    The student generates tokens on-policy from its current distribution $y \sim P_{\text{student}}(\cdot \mid x)$. The teacher evaluates the student's sampled trajectories, and the student optimizes the reverse KL divergence:
    {\boldmath
    \begin{equation}
    \mathcal{L}_{\text{reverse-KL}} = \mathbb{E}_{y \sim P_{\text{student}}} \left[ D_{\text{KL}}\left( P_{\text{student}}(\cdot \mid x, y_{<t}) \,\|\, P_{\text{teacher}}(\cdot \mid x, y_{<t}) \right) \right]
    \end{equation}
    }
    Reverse KL is \emph{zero-forcing} (mode-seeking): the student focuses its limited capacity exclusively on high-likelihood modes of the teacher without hallucinating across non-teacher regions.
\end{enumerate}

\subsection{Expanded Non-Technical Definition (Intuition for Anyone)}
\begin{enumerate}
    \item \textbf{Fixing Student Drift (Exposure Bias)}: If you only teach an apprentice pilot how to fly when the master pilot is steering, the minute the apprentice makes a tiny mistake, they panic and crash. On-policy distillation lets the apprentice fly, catches their specific mistakes in real time, and shows them how the master would handle that exact situation.
    \item \textbf{Mode-Seeking vs. Mode-Covering}: A 70-billion teacher knows 10 different ways to solve a problem. A tiny 2-billion student cannot remember all 10 ways. Forward KL forces the student to try all 10 ways and fail. Reverse KL allows the student to pick the 2 best ways and master them completely!
    \item \textbf{Superhuman Small Models}: Models like TinyLlama, Gemma-2-2B, and LLaMA-3-8B achieve frontier reasoning by distilling hundreds of billions of tokens generated directly from 405B teacher models.
\end{enumerate}

\subsection{Child-Friendly Mental Model (Physical Metaphor)}
Imagine learning to ride a bicycle:
\begin{itemize}
    \item \textbf{Off-Policy Static Distillation}: Watching a video of an Olympic cyclist and copying down the angle of their handlebars. When you get on a real bike and wobble, the video cannot help you!
    \item \textbf{On-Policy Distillation}: You get on the bike. As soon as you wobble left, the coach standing right beside you guides your hands to steer out of the wobble. You learn directly from your own balance adjustments!
\end{itemize}

\section{Core Mathematical Equations and Definitions}

\subsection{Forward vs. Reverse KL Objectives}
{\boldmath
\begin{align}
\mathcal{L}_{\text{fwd}}(p_t, p_s) &= \sum_{v=1}^V p_t(v) \log \left(\frac{p_t(v)}{p_s(v)}\right) \\
\mathcal{L}_{\text{rev}}(p_s, p_t) &= \sum_{v=1}^V p_s(v) \log \left(\frac{p_s(v)}{p_t(v)}\right)
\end{align}
}
\begin{itemize}
    \item \textbf{Formal Definition}: Asymmetric directional information divergences between teacher $p_t$ and student $p_s$ over vocabulary $V$.
    \item \textbf{What It Means}: Forward KL penalizes the student if $p_s(v) \to 0$ when $p_t(v) > 0$. Reverse KL penalizes the student if $p_s(v) > 0$ when $p_t(v) \to 0$.
    \item \textbf{Why It Is Important}: Governs whether the student spreads its probability mass across all teacher options or specializes in crisp modes.
\end{itemize}

\subsection{Jensen-Shannon Divergence Symmetrization}
{\boldmath
\begin{equation}
\mathcal{L}_{\text{JSD}}(p_t, p_s) = \frac{1}{2} D_{\text{KL}}(p_t \,\|\, m) + \frac{1}{2} D_{\text{KL}}(p_s \,\|\, m), \quad m = \frac{p_t + p_s}{2}
\end{equation}
}
\begin{itemize}
    \item \textbf{Formal Definition}: Symmetric bounded divergence metric operating against mixture distribution $m$.
    \item \textbf{What It Means}: Prevents infinite gradient penalties when distributions have disjoint support.
    \item \textbf{Why It Is Important}: Stabilizes training during early phases of distillation when student predictions are uncalibrated.
\end{itemize}

\section{Rigorous Mathematical Analysis Checklist}

\begin{itemize}
    \item \textbf{Notation}: $V$ (vocabulary size), $T$ (sequence length), $p_s, p_t \in \Delta^V$ (next-token probability distributions), $y \in \{0, \dots, V-1\}^T$.
    \item \textbf{Epistemic Status}: Kim \& Rush (EMNLP 2016); Gu et al. (ICLR 2023 - MiniLLM). Standard for autoregressive LLM distillation.
    \item \textbf{Dependency Graph}: Prompt $x \to \text{Sample/Evaluate } y \to \text{Logits at each timestep } t \to \text{Divergence Loss} \to \text{Average over } T$.
    \item \textbf{Dimensions}: Sequence length $T$, vocabulary dimension $V$ (typically $V \in [32000, 128000]$).
    \item \textbf{Regimes}: Edge deployment, high-throughput LLM serving, agentic function calling models.
    \item \textbf{Invariants}: If $p_s \equiv p_t$, all divergences ($\text{KL}_{\text{fwd}}, \text{KL}_{\text{rev}}, \text{JSD}$) vanish to identically $0$.
    \item \textbf{Isomorphisms}: Reverse KL distillation on student rollouts is mathematically isomorphic to PPO policy optimization with teacher log-probability rewards.
\end{itemize}

\section{Theoretical Theorems and Formal Proofs}

\begin{theorem}[Zero-Forcing Nature of Reverse KL Divergence]
Let $p_t, p_s$ be discrete probability distributions over vocabulary $V$. If there exists a token index $v^* \in V$ such that $p_t(v^*) \to 0$, then:
{\boldmath
\begin{equation}
\lim_{p_t(v^*) \to 0} D_{\text{KL}}(p_s \,\|\, p_t) = \infty \quad \text{whenever} \quad p_s(v^*) > 0
\end{equation}
}
thereby forcing any finite-loss student distribution to satisfy $p_s(v) = 0$ for all tokens where $p_t(v) = 0$.
\end{theorem}

\begin{proof}
By definition of Reverse KL divergence:
{\boldmath
\begin{equation}
D_{\text{KL}}(p_s \,\|\, p_t) = \sum_{v=1}^V p_s(v) \log \frac{p_s(v)}{p_t(v)} = \sum_{v \neq v^*} p_s(v) \log \frac{p_s(v)}{p_t(v)} + p_s(v^*) \log \frac{p_s(v^*)}{p_t(v^*)}
\end{equation}
}
Assume $p_s(v^*) = \epsilon > 0$. Taking the limit as $p_t(v^*) \to 0^+$:
{\boldmath
\begin{equation}
\lim_{p_t(v^*) \to 0^+} \log \frac{p_s(v^*)}{p_t(v^*)} = \lim_{p_t(v^*) \to 0^+} \left[ \log \epsilon - \log p_t(v^*) \right] = \infty
\end{equation}
}
Since all terms in the summation $p_s(v) \log \frac{p_s(v)}{p_t(v)} \ge - p_t(v)$ are bounded below, the entire sum diverges to $+\infty$:
{\boldmath
\begin{equation}
\lim_{p_t(v^*) \to 0^+} D_{\text{KL}}(p_s \,\|\, p_t) = +\infty
\end{equation}
}
Under gradient descent, this infinite gradient penalty strictly forces $p_s(v^*) \to 0$. In contrast, for Forward KL:
{\boldmath
\begin{equation}
\lim_{p_t(v^*) \to 0^+} p_t(v^*) \log \frac{p_t(v^*)}{p_s(v^*)} = 0
\end{equation}
}
which incurs zero penalty if the student assigns mass to non-teacher tokens. This proves that Reverse KL is strictly zero-forcing (mode-seeking), preventing hallucinations in small student models.
\end{proof}

\begin{theorem}[Equivalence of On-Policy Reverse KL to Policy Gradient RL]
Minimizing the on-policy sequence reverse KL divergence $\mathcal{L}_{\text{rev}}(\theta) = \mathbb{E}_{y \sim P_\theta}[\log P_\theta(y) - \log P_t(y)]$ is algebraically equivalent to maximizing policy gradient reinforcement learning with reward $R(y) = \log P_t(y) - \log P_\theta(y)$.
\end{theorem}

\begin{proof}
Let $J(\theta) = - \mathcal{L}_{\text{rev}}(\theta) = \mathbb{E}_{y \sim P_\theta}[\log P_t(y) - \log P_\theta(y)]$.
Using the policy gradient theorem and log-derivative identity $\nabla_\theta P_\theta(y) = P_\theta(y) \nabla_\theta \log P_\theta(y)$:
{\boldmath
\begin{align}
\nabla_\theta J(\theta) &= \nabla_\theta \sum_{y} P_\theta(y) \left[ \log P_t(y) - \log P_\theta(y) \right] \\
&= \sum_{y} \left( \nabla_\theta P_\theta(y) \left[ \log P_t(y) - \log P_\theta(y) \right] - P_\theta(y) \frac{\nabla_\theta P_\theta(y)}{P_\theta(y)} \right) \\
&= \sum_{y} P_\theta(y) \nabla_\theta \log P_\theta(y) \left[ \log P_t(y) - \log P_\theta(y) \right] - \sum_{y} \nabla_\theta P_\theta(y)
\end{align}
}
Since $\sum_y P_\theta(y) = 1$, the second term $\nabla_\theta \sum_y P_\theta(y) = \nabla_\theta(1) = 0$.
Therefore:
{\boldmath
\begin{equation}
\nabla_\theta J(\theta) = \mathbb{E}_{y \sim P_\theta} \left[ \nabla_\theta \log P_\theta(y) \cdot R(y) \right], \quad R(y) = \log P_t(y) - \log P_\theta(y)
\end{equation}
}
This matches the exact formulation of the REINFORCE policy gradient objective with reward $R(y)$, proving that on-policy distillation is isomorphic to reinforcement learning where the teacher acts as the environment reward signal.
\end{proof}

\section{Foundational Architectural Questions and Epistemic Meta-Questions}

\subsection{Architectural Question 1: Mitigating Off-Distribution Degeneration}
\textbf{Question:} Why does sequence-level distillation often yield higher accuracy than token-level forward KL on static data?\\
\textbf{Answer:} Static human-written transcripts often contain high-entropy lexical variations that the student's limited parameter budget cannot capture. When the teacher generates sequences autoregressively, it picks consistent high-probability paths that eliminate unnatural variance, producing a cleaner, more learnable data manifold for the student.

\textbf{Meta-Question:} How should length penalties and temperature be tuned during teacher sequence generation?\\
\textbf{Answer:} Teacher rollouts typically use moderate temperature ($T \approx 0.7$) with nucleus sampling ($p = 0.9$) or greedy decoding ($T = 0$). Greedy decoding produces the highest deterministic precision for reasoning and code distillation, while nucleus sampling produces diverse conversational corpora.

\subsection{Architectural Question 2: Serving Costs of Distillation}
\textbf{Question:} How is teacher inference parallelized when generating billions of distillation tokens?\\
\textbf{Answer:} Sequence-level distillation separates teacher rollout generation from student training into an asynchronous offline pipeline:
1. Batches of prompts are sent to teacher serving clusters running TensorRT-LLM or vLLM with FP8/FP4 acceleration.
2. Complete teacher responses and token log-probabilities are serialized to high-throughput object stores.
3. The student trains via high-speed multi-node data-parallel SFT, decoupling training throughput from teacher inference latency.

\textbf{Meta-Question:} Can student models surpass teacher capabilities through distillation?\\
\textbf{Answer:} In standard sequence distillation, students represent an upper-bounded compression of the teacher. However, when combined with verifiable reinforcement learning (RLVR, #214) or rejection sampling (Best-of-N, #216), student models frequently surpass teacher performance on mathematical and algorithmic benchmarks.

\section{Algorithmic Implementation Specification}

\begin{algorithm}[H]
\caption{Autoregressive Sequence Distillation Stepwise Loss}
\begin{algorithmic}[1]
\Require Student logits $\mathbf{Z}^s \in \mathbb{R}^{T \times V}$, teacher logits $\mathbf{Z}^t \in \mathbb{R}^{T \times V}$, tokens $\mathbf{y} \in \{0, \dots, V-1\}^T$, mode $\in \{\text{forward\_kl}, \text{reverse\_kl}, \text{sequence\_sft}, \text{jsd}\}$, temperature $T_0$
\Ensure Token losses $\mathbf{L} \in \mathbb{R}^T$, mean loss $\bar{L}$, student entropies $\mathbf{H}^s \in \mathbb{R}^T$
\State Allocate $\mathbf{L} \in \mathbb{R}^T, \quad \mathbf{H}^s \in \mathbb{R}^T$
\For{$t = 0$ \textbf{to} $T-1$}
    \State $\mathbf{p}^s \gets \operatorname{Softmax}(\mathbf{Z}^s[t] / T_0), \quad \mathbf{p}^t \gets \operatorname{Softmax}(\mathbf{Z}^t / T_0)$
    \State $\mathbf{H}^s[t] \gets - \sum_{v=0}^{V-1} \mathbf{p}^s[v] \ln(\max(10^{-12}, \mathbf{p}^s[v]))$
    \If{$\text{mode} == \text{"reverse\_kl"}$}
        \State $\mathbf{L}[t] \gets \sum_{v=0}^{V-1} \mathbf{p}^s[v] \ln\left( \frac{\max(10^{-12}, \mathbf{p}^s[v])}{\max(10^{-12}, \mathbf{p}^t[v])} \right)$
    \ElsIf{$\text{mode} == \text{"sequence\_sft"}$}
        \State $\mathbf{L}[t] \gets - \ln(\max(10^{-12}, \mathbf{p}^s[\mathbf{y}[t]])))$
    \ElsIf{$\text{mode} == \text{"jsd"}$}
        \State $\mathbf{m} \gets 0.5 \cdot (\mathbf{p}^t + \mathbf{p}^s)$
        \State $\mathbf{L}[t] \gets 0.5 \left( D_{\text{KL}}(\mathbf{p}^t \,\|\, \mathbf{m}) + D_{\text{KL}}(\mathbf{p}^s \,\|\, \mathbf{m}) \right)$
    \Else
        \State $\mathbf{L}[t] \gets \sum_{v=0}^{V-1} \mathbf{p}^t[v] \ln\left( \frac{\max(10^{-12}, \mathbf{p}^t[v])}{\max(10^{-12}, \mathbf{p}^s[v])} \right)$
    \EndIf
\EndFor
\State $\bar{L} \gets \frac{1}{T} \sum_{t=0}^{T-1} \mathbf{L}[t]$
\State \Return $(\mathbf{L}, \bar{L}, \mathbf{H}^s)$
\end{algorithmic}
\end{algorithm}

\section{Big-$\Theta$ Complexity Analysis}

\begin{itemize}
    \item \textbf{Softmax Distribution Time}: $\Theta(T \cdot V)$ operations across $T$ tokens.
    \item \textbf{Divergence Loss Evaluation}: $\Theta(T \cdot V)$ for Forward/Reverse KL and JSD; $\Theta(T)$ for Sequence SFT.
    \item \textbf{Shannon Entropy Computation}: $\Theta(T \cdot V)$ operations.
    \item \textbf{Total Time Complexity}: $\Theta(T \cdot V)$ operations per training sequence.
    \item \textbf{Memory Storage Overhead}: $\Theta(T \cdot V)$ floats for intermediate logit buffers.
\end{itemize}

\section{Single Canonical Repository Reference}

The complete deterministic scalar implementation, verification suite, and unit test benchmarks are published at:
\begin{center}
\url{https://github.com/Chief-Strategist-J/llm-obs-infra}
\end{center}

\end{document}
